% !TeX document-id = {fc72019d-d3ba-4743-9e76-a08bcafbe161}
% !TeX TXS-program:compile = txs:///pdflatex/[-shell-escape]

\documentclass[11pt]{beamer}
\usepackage[utf8]{inputenc}
\usepackage[T1,T2A]{fontenc}
\usepackage[russian]{babel}
\usepackage{color}
\usepackage{listings}
\usepackage{calc}
\usepackage{graphicx}
\usepackage{epstopdf}
\usepackage{hyperref}
\hypersetup{unicode,colorlinks}
\usepackage{csquotes}
\usepackage{upquote}
\usepackage{cprotect}
\usetheme[progressbar=head,numbering=fraction,block=fill]{metropolis}
\usepackage{dejavu}
\usepackage{etoolbox}
\usepackage{bm}
\usepackage{minted}
%\usepackage[defaultsans]{droidsans}
%\usepackage[defaultmono]{droidmono}
%\makeatletter
%\input droid/t1fdm.fd
%\makeatother

\makeatletter
\DeclareRobustCommand*{\bfseries}{%
    \not@math@alphabet\bfseries\mathbf
    \fontseries\bfdefault\selectfont
    \boldmath
}
\makeatother

\makeatletter
\chardef\straightquote@code=\catcode`'
\chardef\backquote@code=\catcode``
\catcode`'=\active \catcode``=\active
\patchcmd{\@noligs}
{\textasciigrave}
{\fixedtextasciigrave}
{}{}
\newcommand{\fixedtextasciigrave}{%
    \makebox[.5em]{\fontencoding{TS1}\fontfamily{fvs}\selectfont\textasciigrave}% Vera Sans
}
\catcode\lq\'=\straightquote@code
\catcode\lq\`=\backquote@code
\makeatletter

\lstset{
    language=Haskell,
    inputencoding=utf8,
    extendedchars=true,
    breaklines=true,
    escapeinside=!!,
    tabsize=4,
    breakatwhitespace=true,
    keepspaces=true,
    basicstyle=\ttfamily,
    commentstyle=\color{gray}\itshape,
    keywordstyle=
    % keywordstyle=\color{blue}, % TODO https://tex.stackexchange.com/questions/415777/avoid-highlighting-keywords-following-certain-words-in-listings
    % addkeywords={Type}
    % deletekeywords={Monad, String, return, array, show},
    % otherkeywords={::, ->}
}
\lstset{
    literate={а}{{\selectfont\char224}}1
    {б}{{\selectfont\char225}}1
    {в}{{\selectfont\char226}}1
    {г}{{\selectfont\char227}}1
    {д}{{\selectfont\char228}}1
    {е}{{\selectfont\char229}}1
    {ё}{{\"e}}1
    {ж}{{\selectfont\char230}}1
    {з}{{\selectfont\char231}}1
    {и}{{\selectfont\char232}}1
    {й}{{\selectfont\char233}}1
    {к}{{\selectfont\char234}}1
    {л}{{\selectfont\char235}}1
    {м}{{\selectfont\char236}}1
    {н}{{\selectfont\char237}}1
    {о}{{\selectfont\char238}}1
    {п}{{\selectfont\char239}}1
    {р}{{\selectfont\char240}}1
    {с}{{\selectfont\char241}}1
    {т}{{\selectfont\char242}}1
    {у}{{\selectfont\char243}}1
    {ф}{{\selectfont\char244}}1
    {х}{{\selectfont\char245}}1
    {ц}{{\selectfont\char246}}1
    {ч}{{\selectfont\char247}}1
    {ш}{{\selectfont\char248}}1
    {щ}{{\selectfont\char249}}1
    {ъ}{{\selectfont\char250}}1
    {ы}{{\selectfont\char251}}1
    {ь}{{\selectfont\char252}}1
    {э}{{\selectfont\char253}}1
    {ю}{{\selectfont\char254}}1
    {я}{{\selectfont\char255}}1
    {А}{{\selectfont\char192}}1
    {Б}{{\selectfont\char193}}1
    {В}{{\selectfont\char194}}1
    {Г}{{\selectfont\char195}}1
    {Д}{{\selectfont\char196}}1
    {Е}{{\selectfont\char197}}1
    {Ё}{{\"E}}1
    {Ж}{{\selectfont\char198}}1
    {З}{{\selectfont\char199}}1
    {И}{{\selectfont\char200}}1
    {Й}{{\selectfont\char201}}1
    {К}{{\selectfont\char202}}1
    {Л}{{\selectfont\char203}}1
    {М}{{\selectfont\char204}}1
    {Н}{{\selectfont\char205}}1
    {О}{{\selectfont\char206}}1
    {П}{{\selectfont\char207}}1
    {Р}{{\selectfont\char208}}1
    {С}{{\selectfont\char209}}1
    {Т}{{\selectfont\char210}}1
    {У}{{\selectfont\char211}}1
    {Ф}{{\selectfont\char212}}1
    {Х}{{\selectfont\char213}}1
    {Ц}{{\selectfont\char214}}1
    {Ч}{{\selectfont\char215}}1
    {Ш}{{\selectfont\char216}}1
    {Щ}{{\selectfont\char217}}1
    {Ъ}{{\selectfont\char218}}1
    {Ы}{{\selectfont\char219}}1
    {Ь}{{\selectfont\char220}}1
    {Э}{{\selectfont\char221}}1
    {Ю}{{\selectfont\char222}}1
    {Я}{{\selectfont\char223}}1
    {`}{{{\fixedtextasciigrave}}}1
}

% \usemintedstyle{solarized-light}
\newminted[haskell]{haskell}{
    escapeinside=``,
    beameroverlays=true,
    texcomments=true,
    autogobble=true,
    fontsize=\small
}
\newmintinline[haskinline]{haskell}{
    escapeinside=``,
    beameroverlays=true,
    texcomments=true
}

\author{Алексей Романов}
\subtitle{Функциональное программирование на Haskell}
%\logo{}
\institute{МИЭТ}
\subject{Функциональное программирование на Haskell}
%\setbeamercovered{transparent}
%\setbeamertemplate{navigation symbols}{}


\title{Лекция 17: что дальше?}

\begin{document}
\begin{frame}[plain]
  \maketitle
\end{frame}

\begin{frame}[fragile]
  \frametitle{Веб-интерфейсы}
  \small
  Что происходит при нажатии кнопки?
  \begin{itemize}
    \item \href{https://haskell-miso.org/}{Miso}: код на Haskell компилируется в JavaScript или WebAssembly. В браузере нажатие обновляет модель и представление.
    \item \href{https://data-star.dev/}{Datastar} и \hackage{datastar-hs}: сервер на Haskell может ответить на нажатие фрагментом HTML. Браузер обновляет часть страницы. Пример: \href{https://carlo-hamalainen.net/2026/08/06/ghc-vis-ds/}{ghc-vis-ds} из лекции 9.
    \item \href{https://elm-lang.org/}{Elm} "--- отдельный функциональный язык для веб-интерфейсов. В \href{https://guide.elm-lang.org/architecture/}{The Elm Architecture} нажатие порождает сообщение. Функция \texttt{update} вычисляет новую модель, а \texttt{view} описывает, что увидит пользователь.
  \end{itemize}
\end{frame}

\begin{frame}
  \frametitle{Настольные графические интерфейсы}
  \small
  \begin{itemize}
    \item \href{https://github.com/fjvallarino/monomer}{Monomer}: GUI для Windows, macOS и Linux. Как и в Elm, события обновляют модель интерфейса.
    \item \href{https://github.com/haskell-gi/haskell-gi}{gi-gtk}: виджеты \href{https://www.gtk.org/}{GTK} и обработчики событий в \texttt{IO}.
          \begin{itemize}
            \item Для GTK 3 есть \href{https://github.com/owickstrom/gi-gtk-declarative}{gi-gtk-declarative}: интерфейс описывается деревом виджетов.
          \end{itemize}
    \item \href{https://github.com/haskell-game/dear-imgui.hs}{dear-imgui.hs} "--- интерфейс в немедленном режиме (\emph{immediate mode}). В каждом кадре вызываем \texttt{button} и сразу узнаём, нажата ли кнопка.
    \item \href{https://github.com/jtdaugherty/brick}{brick}: интерактивные интерфейсы в терминале.
    \item \href{https://bradrn.com/posts/hs-gui-libs.html}{The Big List of Haskell GUI Libraries} (2022): обзор библиотек и сравнение подходов.
  \end{itemize}
\end{frame}

\begin{frame}[fragile]
  \frametitle{Библиотеки и инструменты Haskell}
  \begin{itemize}
    \item Сборка проектов и управление зависимостями: \href{https://cabal.readthedocs.io/en/stable/}{cabal} и \href{https://docs.haskellstack.org/en/stable/}{stack}.
    \item \href{https://haskell-language-server.readthedocs.io/en/latest/features.html}{Haskell Language Server}: показывает типы при наведении в редакторе и предлагает, чем заполнить типизированные дыры.
    \item Сериализация значений Haskell: \hackage{aeson} для JSON, \hackage{binary} для двоичного представления.
    \item \href{https://hackage.haskell.org/package/criterion/docs/Criterion-Main.html}{criterion}: измерение производительности. Глубину вычисления результата задают \haskinline|whnf| и \haskinline|nf|.
    \item Линзы (\hackage{lens}): обобщение геттеров и сеттеров.
    \item Альтернативы \haskinline|show| для больших структур: \hackage{pretty-show} и \hackage{prettyprinter}.
  \end{itemize}
\end{frame}

\begin{frame}
  \frametitle{Серверная часть веб-приложения}
  \begin{itemize}
    \item \href{https://github.com/scotty-web/scotty/tree/master/examples}{Scotty}: маршруты, обработка HTTP-запросов и ответы в JSON.
    \item \href{https://www.yesodweb.com/book}{Yesod}: фреймворк с типизированными маршрутами, HTML-шаблонами и формами. В книге разобран пример блога.
    \item \href{https://aosabook.org/en/posa/warp.html}{Warp}: как устроен HTTP-сервер, на котором можно запускать приложения на Scotty и Yesod.
    \item HTTP-клиент \href{https://hackage.haskell.org/package/http-conduit/docs/Network-HTTP-Simple.html}{\texttt{Network.HTTP.Simple}}: \href{https://hackage.haskell.org/package/http-conduit/docs/Network-HTTP-Simple.html\#v:httpJSON}{\texttt{httpJSON}} отправляет запрос другому сервису и разбирает JSON в ответе.
    \item \href{https://www.yesodweb.com/book/persistent}{persistent}: работа с SQLite, PostgreSQL и MySQL через записи Haskell.
    \item \href{https://github.com/bitemyapp/esqueleto}{esqueleto}: типобезопасные SQL-запросы для persistent, в том числе с \texttt{JOIN}.
  \end{itemize}
\end{frame}

\begin{frame}[fragile]
  \frametitle{Нерассмотренные части Haskell: язык}
  \small
  \begin{itemize}
    \item \href{https://downloads.haskell.org/ghc/latest/docs/users_guide/exts/generics.html}{\haskinline{deriving Generic}} даёт общее представление конструкторов и полей. По нему можно задать правило вывода экземпляров для своего класса. Так делают \hackage{aeson} и \hackage{binary}.
    \item \href{https://downloads.haskell.org/ghc/latest/docs/users_guide/exts/type_families.html}{Семейства типов}.
    \item \href{https://downloads.haskell.org/ghc/latest/docs/users_guide/exts/rank_polymorphism.html}{Типы высших рангов} с \haskinline|forall| \emph{внутри} типа (пример: \href{https://hackage.haskell.org/package/base/docs/Control-Monad-ST.html#v:runST}{\haskinline{runST}}).
    \item \href{https://downloads.haskell.org/ghc/latest/docs/users_guide/exts/existential_quantification.html}{Экзистенциальные типы}.
    \item \href{https://downloads.haskell.org/ghc/latest/docs/users_guide/exts/pattern_synonyms.html}{Синонимы образцов}.
    \item \href{https://downloads.haskell.org/ghc/latest/docs/users_guide/exts/template_haskell.html}{Template Haskell}: генерация кода при компиляции (см. примеры в QuickCheck и Hedgehog).
    \item Вызов кода на других языках: \href{https://downloads.haskell.org/ghc/latest/docs/users_guide/exts/ffi.html}{FFI} для \href{https://www.open-std.org/jtc1/sc22/wg14/}{C}, \hackage{inline-c-cpp} для \href{https://www.open-std.org/jtc1/sc22/wg21/}{C++} (используется в \href{https://github.com/haskell-game/dear-imgui.hs}{dear-imgui.hs}), \href{https://github.com/harmont-dev/hsrs}{hsrs} для \href{https://rust-lang.org/}{Rust}.
  \end{itemize}
\end{frame}

\begin{frame}[fragile]
  \frametitle{Подходы к эффектам}
  \begin{itemize}
    \item Многие монады (и аппликативные функторы) можно рассматривать как побочные эффекты.
          \begin{itemize}
            \item \haskinline|IO|: понятно.
            \item \haskinline|State|: чтение и запись в память.
            \item \haskinline|[]|: недетерминированность
            \item \ldots
          \end{itemize}
          \pause
    \item Как их комбинировать?
    \item Композиция монад не обязательно монада.
          \pause
    \item Трансформеры монад: \hackage{transformers} и \hackage{mtl}
    \item \href{https://www.microsoft.com/en-us/research/wp-content/uploads/2016/08/algeff-tr-2016-v2.pdf}{Алгебраические эффекты}
    \item \href{https://haskellforall.com/2012/06/you-could-have-invented-free-monads}{Свободные (free)} и \href{https://okmij.org/ftp/Haskell/extensible/more.pdf}{\enquote{более свободные} (freer)} монады
    \item Небольшой обзор: \href{https://ocharles.org.uk/posts/2016-01-26-transformers-free-monads-mtl-laws.html}{Monad transformers, free monads, mtl, laws and a new approach}
  \end{itemize}
\end{frame}

\begin{frame}[fragile]
  \frametitle{Будущее (и недавнее прошлое) Haskell}
  \begin{itemize}
    \item \href{https://downloads.haskell.org/ghc/latest/docs/users_guide/exts/linear_types.html}{Линейные типы} (GHC 9.0+): значения, которые можно использовать только один раз.
          \begin{itemize}
            \item Меньше сборки мусора.
            \item Можно использовать изменяемые структуры данных.
            \item Безопасное изменение состояния типа (например, тип файла может включать, открыт ли он).
          \end{itemize}
          \pause
    \item \href{https://serokell.io/blog/2018/12/17/why-dependent-haskell}{Зависимые типы}: типы, параметризованные значениями.
          \begin{itemize}
            \item Стандартный пример: вектор с размером в типе.
            \item Библиотека \hackage{singletons} позволяет эмулировать многие применения зависимых типов.
          \end{itemize}
  \end{itemize}
\end{frame}

\begin{frame}[fragile]
  \frametitle{Альтернативы Haskell}
  \small
  \begin{itemize}
    \setlength{\itemsep}{1pt}
    \item \href{https://ocaml.org/}{OCaml}: похож на \href{https://www.haskell.org/}{Haskell} без ленивости и классов типов.
    \item \href{https://www.erlang.org/}{Erlang}: функциональный язык с моделью акторов (процессов, взаимодействующих через сообщения).
    \item Семейство Lisp: \href{https://lisp-lang.org/}{Common Lisp} и \href{https://racket-lang.org/}{Racket}.
          \pause
    \item Для JVM: \href{https://www.scala-lang.org/}{Scala}, \href{https://clojure.org/}{Clojure}.
    \item Для .Net: \href{https://fsharp.org/}{F\#} (почти тот же \href{https://ocaml.org/}{OCaml}).
          \pause
    \item \href{https://rust-lang.org/}{Rust}: альтернатива \href{https://www.open-std.org/jtc1/sc22/wg21/}{C++} без сборщика мусора; типы отслеживают время жизни значений.
          \pause
    \item Языки с зависимыми типами: \href{https://agda.readthedocs.io/}{Agda}, \href{https://idris-lang.org/}{Idris}, \href{https://rocq-prover.org/}{Rocq}.
          \pause
    \item Языки спецификаций: \href{https://lamport.azurewebsites.net/tla/tla.html}{TLA+}, \href{https://alloytools.org/}{Alloy}.
  \end{itemize}
\end{frame}

\begin{frame}[fragile]
  \frametitle{Области математики, связанные с ФП}
  \begin{itemize}
    \item Лямбда-исчисление: то, с чего всё началось.
          \begin{itemize}
            \item Питер Селингер, \href{https://arxiv.org/abs/0804.3434}{Lecture Notes on the Lambda Calculus}: от бестипового исчисления до полиморфизма и вывода типов.
          \end{itemize}
          \pause
    \item Математическая логика в целом
          \begin{itemize}
            \item Логики высших порядков.
            \item Изоморфизм Карри-Ховарда.
            \item Формализация математики.
            \item И автоматический поиск доказательств.
            \item \href{https://agda.readthedocs.io/}{Agda}, \href{https://idris-lang.org/}{Idris}, \href{https://rocq-prover.org/}{Rocq (бывший Coq)} и \href{https://lean-lang.org/}{Lean} сюда же!
          \end{itemize}
          \pause
    \item Теория категорий: источник понятий функтора, монады, естественного преобразования.
          \begin{itemize}
            \item Bartosz Milewski, \href{https://github.com/hmemcpy/milewski-ctfp-pdf}{Category Theory for Programmers}.
          \end{itemize}
  \end{itemize}
\end{frame}


\begin{frame}[fragile]
  \frametitle{Дополнительное чтение}
  \begin{itemize}
    \item \href{https://sdiehl.github.io/wiwinwlh/}{What I Wish I Knew When Learning Haskell}
    \item \href{https://lexi-lambda.github.io/blog/2018/02/10/an-opinionated-guide-to-haskell-in-2018/}{An opinionated guide to Haskell in 2018}
    \item \href{https://realworldocaml.org/}{Real World OCaml: Functional Programming for the Masses} (бесплатная книга)
    \item \href{https://wiki.portal.chalmers.se/agda/pmwiki.php?n=Main.Documentation}{Много ссылок по Agda}
    \item \href{https://idris2.readthedocs.io/en/latest/tutorial/index.html}{Введение в Idris}
    \item \href{https://doc.rust-lang.org/book/}{The Rust Programming Language} (бесплатная книга)
  \end{itemize}
\end{frame}

\end{document}
